\documentclass[11pt]{article}

\usepackage[margin=1in]{geometry}
\usepackage{amsmath,amssymb,mathtools}
\usepackage{booktabs,array,longtable}
\usepackage{enumitem}
\usepackage{xcolor}
\usepackage{hyperref}
\usepackage{listings}

\hypersetup{
  colorlinks=true,
  linkcolor=blue,
  urlcolor=blue,
  citecolor=blue
}

\lstset{
  basicstyle=\ttfamily\small,
  columns=fullflexible,
  breaklines=true,
  frame=single,
  showstringspaces=false
}

\newcommand{\Hcal}{\mathcal{H}}
\newcommand{\Scal}{\mathcal{S}}
\newcommand{\Acal}{\mathcal{A}}
\newcommand{\Prb}{\mathbb{P}}
\setlist[description]{style=nextline,leftmargin=0pt}
\emergencystretch=2em

\title{Aggregate Documentation Report for \texttt{utils\_jin}}
\author{STAT 221 Bayesian Poker Project}
\date{May 1, 2026}

\begin{document}
\maketitle

\tableofcontents
\newpage

\section{Purpose and Scope}

The \texttt{utils\_jin/} package is the first complete implementation layer for
running pre-flop Bayesian range inference and preparing the project for EM. It
does not replace the original \texttt{utils/} folder. Instead, it wraps and
extends the existing parser so that the model code has clean, compact objects:

\[
\begin{aligned}
\text{PHH sessions}
&\rightarrow \text{parsed hands}
\rightarrow \text{decision records} \\
&\rightarrow \text{state/action encodings}
\rightarrow \text{pre-flop filtering} \\
&\rightarrow \text{EM soft counts}.
\end{aligned}
\]

The package currently focuses on pre-flop inference. It includes post-flop
state keys and records because the data extractor already sees all streets, but
the implemented filter and EM starter are pre-flop only. This is intentional:
pre-flop filtering is the first working model, and post-flop strength filtering
should plug into the same record/state infrastructure later.

\section{High-Level Architecture}

\subsection{Main Data Flow}

The current intended data flow is:

\begin{enumerate}[leftmargin=*]
  \item \texttt{compat.py} makes the legacy \texttt{utils/parse.py} parser importable.
  \item \texttt{records.py} parses hands and produces one \texttt{DecisionRecord}
  per player action.
  \item \texttt{action\_encoder.py} recomputes model-facing action buckets.
  \item \texttt{state\_encoder.py} maps rich parser states to small discrete
  \texttt{StateKey} values.
  \item \texttt{cards.py} provides 1326-combo and 169-class card utilities.
  \item \texttt{gto\_prior.py} provides a simple heuristic
  \(p_{\mathrm{GTO}}(a \mid h,s)\).
  \item \texttt{preflop\_filter.py} updates a range
  \(R_t(h)\) using Bayesian filtering.
  \item \texttt{em\_preflop.py} turns filtering posteriors into EM
  responsibilities and soft counts.
\end{enumerate}

\subsection{Core Mathematical Objects}

\begin{center}
\begin{tabular}{lll}
\toprule
\textbf{Object} & \textbf{Meaning} & \textbf{Implemented In} \\
\midrule
\(h \in \Hcal\) & 169 pre-flop hand class & \texttt{cards.py} \\
\(s \in \Scal\) & compact public state key & \texttt{state\_encoder.py} \\
\(a \in \Acal\) & action bucket & \texttt{action\_encoder.py} \\
\(p_{\mathrm{GTO}}(a\mid h,s)\) & heuristic pre-flop action prior & \texttt{gto\_prior.py} \\
\(R_t(h)\) & filtered range after action \(t\) & \texttt{preflop\_filter.py} \\
\(q_m(h)\) & EM responsibility for player-hand \(m\) & \texttt{em\_preflop.py} \\
\bottomrule
\end{tabular}
\end{center}

\subsection{Bayesian Filtering Equation}

The pre-flop filter performs the recursive update:

\[
R_t(h)
=
\frac{
R_{t-1}(h)\,\phi_{h,s_t}[a_t]
}{
\sum_{h' \in \Hcal} R_{t-1}(h')\,\phi_{h',s_t}[a_t]
}.
\]

In the current implementation,
\[
\phi_{h,s}[a] = p_{\mathrm{GTO}}(a \mid h,s),
\]
where \(p_{\mathrm{GTO}}\) is a transparent heuristic, not a solver output.
The interface is designed so that a real solver table can later replace this
heuristic without rewriting the filter.

\subsection{EM Preparation}

The implemented EM starter uses the filter as the E-step engine. For each
player-hand group \(m\), it computes:

\[
q_m(h)
=
\Prb(V_m=h \mid a^m_{1:T_m}, s^m_{1:T_m}).
\]

Then it forms soft counts:

\[
N_{h,s,a}
=
\sum_m
\sum_{t \in \mathrm{pre}}
q_m(h)\,
\mathbf{1}\{s_t^m=s,\ a_t^m=a\}.
\]

Finally, the M-step helper computes MAP-style action likelihoods:

\[
\phi_{h,s,a}
\propto
N_{h,s,a}
+ \lambda p_{\mathrm{GTO}}(a\mid h,s)
+ \epsilon.
\]

\section{Package Inventory}

\begin{center}
\begin{tabular}{lrp{8.3cm}}
\toprule
\textbf{File} & \textbf{Lines} & \textbf{Primary Role} \\
\midrule
\texttt{\_\_init\_\_.py} & 13 & Package description and namespace marker. \\
\texttt{compat.py} & 48 & Compatibility wrapper for importing legacy parser modules. \\
\texttt{cards.py} & 207 & Card normalization, 1326 combos, 169 classes, priors. \\
\texttt{action\_encoder.py} & 153 & Model-facing action buckets using pre-action context. \\
\texttt{state\_encoder.py} & 203 & Compact public state keys for filtering and EM. \\
\texttt{records.py} & 224 & Flat decision-record extraction from parsed hands. \\
\texttt{gto\_prior.py} & 218 & Simple GTO-inspired pre-flop action prior. \\
\texttt{preflop\_filter.py} & 124 & Bayesian range filter over 169 classes. \\
\texttt{em\_preflop.py} & 139 & E-step, soft-count, and M-step helpers. \\
\texttt{README.md} & 47 & Short usage guide. \\
\bottomrule
\end{tabular}
\end{center}

\newpage
\section{\texttt{\_\_init\_\_.py}}

\subsection{Role}

\texttt{\_\_init\_\_.py} turns \texttt{utils\_jin/} into a Python package and
documents the purpose of the package. It intentionally contains no runtime
logic. Its docstring states that the package provides:

\begin{itemize}[leftmargin=*]
  \item card and combo helpers;
  \item compact public state encoders;
  \item a simple GTO-inspired pre-flop action prior;
  \item a pre-flop Bayesian range filter;
  \item EM starter functions.
\end{itemize}

\subsection{Why It Exists}

Keeping \texttt{utils\_jin/} as a package means downstream code can import:

\begin{lstlisting}[language=Python]
from utils_jin.records import extract_session_records
from utils_jin.preflop_filter import filter_preflop_records
\end{lstlisting}

This is cleaner than relying on notebook-relative imports.

\section{\texttt{compat.py}}

\subsection{Problem Solved}

The existing parser in \texttt{utils/parse.py} imports \texttt{hand\_map} and
\texttt{action\_map} as top-level modules:

\begin{lstlisting}[language=Python]
from hand_map import poker_hand_mapper
from action_map import classify
\end{lstlisting}

That works when a notebook runs from inside \texttt{utils/}, but it is fragile
when importing from the repository root. \texttt{compat.py} solves this without
editing legacy code by temporarily adding the old \texttt{utils/} directory to
\texttt{sys.path}.

\subsection{Functions}

\begin{description}
  \item[\texttt{repo\_root()}] Returns the repository root based on the location
  of \texttt{utils\_jin/compat.py}.
  \item[\texttt{ensure\_legacy\_utils\_path()}] Adds the legacy \texttt{utils/}
  path to \texttt{sys.path} if needed.
  \item[\texttt{load\_legacy\_session\_class()}] Imports and returns the original
  \texttt{Session} class from \texttt{utils/parse.py}.
  \item[\texttt{load\_legacy\_hand\_mapper()}] Imports and returns the original
  \texttt{poker\_hand\_mapper}.
\end{description}

\subsection{Design Note}

This file is a reconciliation layer. It avoids touching the old parser while
making the new package robust. If the legacy parser is later converted into a
proper package, this compatibility module can be simplified or removed.

\section{\texttt{cards.py}}

\subsection{Role}

\texttt{cards.py} defines the card universe and the mapping between exact
two-card combinations and 169 pre-flop classes. It supports both the modeling
view and the poker view:

\begin{itemize}[leftmargin=*]
  \item exact combo space: \(\binom{52}{2}=1326\);
  \item equivalence-class space: 169 classes such as \texttt{AA}, \texttt{AKs},
  and \texttt{AKo}.
\end{itemize}

The current pre-flop filter tracks 169 classes. Exact combos are still needed
for future post-flop blocker-aware mapping.

\subsection{Constants}

\begin{center}
\begin{tabular}{ll}
\toprule
\textbf{Name} & \textbf{Meaning} \\
\midrule
\texttt{RANKS\_LOW\_TO\_HIGH} & Rank order from 2 through A. \\
\texttt{RANKS\_HIGH\_TO\_LOW} & Rank order from A through 2. \\
\texttt{SUITS} & Suit order \texttt{s}, \texttt{h}, \texttt{d}, \texttt{c}. \\
\texttt{RANK\_TO\_VALUE} & Rank-to-integer map, A = 14. \\
\texttt{VALUE\_TO\_RANK} & Inverse map. \\
\bottomrule
\end{tabular}
\end{center}

\subsection{Important Functions}

\begin{description}
  \item[\texttt{normalize\_card(card)}] Standardizes cards to rank-uppercase and
  suit-lowercase, e.g. \texttt{AS} becomes \texttt{As}.
  \item[\texttt{split\_cards(cards)}] Converts compact strings such as
  \texttt{AsKd} into \texttt{["As", "Kd"]}.
  \item[\texttt{full\_deck()}] Returns all 52 cards in stable order.
  \item[\texttt{all\_combos()}] Returns all 1326 exact two-card combos.
  \item[\texttt{combo\_to\_class(combo)}] Maps exact cards to a 169 class.
  \item[\texttt{all\_169\_classes()}] Returns all 169 pre-flop classes.
  \item[\texttt{class\_multiplicity(hand\_class)}] Returns 6 for pairs, 4 for
  suited non-pairs, and 12 for offsuit non-pairs.
  \item[\texttt{initial\_class\_prior()}] Returns the class prior
  \(C(h)/1326\).
  \item[\texttt{legal\_combos(dead\_cards)}] Removes exact combos blocked by
  board or known cards.
  \item[\texttt{aggregate\_combo\_probs(combo\_probs)}] Aggregates exact-combo
  probabilities into 169-class probabilities.
\end{description}

\subsection{Examples}

\begin{lstlisting}[language=Python]
from utils_jin.cards import combo_to_class, all_combos, initial_class_prior

combo_to_class("AsKs")  # "AKs"
combo_to_class("AsKd")  # "AKo"
combo_to_class("AhAd")  # "AA"

len(all_combos())       # 1326
sum(initial_class_prior().values())  # 1.0
\end{lstlisting}

\subsection{Model Connection}

The initial range \(R_0\) is:

\[
R_0(h)=\frac{C(h)}{1326},
\qquad
C(h)=
\begin{cases}
6 & \text{pair},\\
4 & \text{suited non-pair},\\
12 & \text{offsuit non-pair}.
\end{cases}
\]

This is implemented by \texttt{initial\_class\_prior()}.

\section{\texttt{action\_encoder.py}}

\subsection{Role}

\texttt{action\_encoder.py} converts legacy parser action tuples into clean
model-facing action buckets. The legacy parser already stores an action bucket,
but it computes bet/raise size using the pot after the action. This script
recomputes the bucket from the pre-action state, which is more appropriate for
modeling.

\subsection{Action Buckets}

\begin{center}
\begin{tabular}{cll}
\toprule
\textbf{ID} & \textbf{Constant} & \textbf{Label} \\
\midrule
0 & \texttt{FOLD} & \texttt{fold} \\
1 & \texttt{CHECK\_CALL} & \texttt{check\_call} \\
2 & \texttt{SMALL\_BET\_RAISE} & \texttt{small\_bet\_raise} \\
3 & \texttt{MEDIUM\_BET\_RAISE} & \texttt{medium\_bet\_raise} \\
4 & \texttt{LARGE\_BET\_RAISE} & \texttt{large\_bet\_raise} \\
\bottomrule
\end{tabular}
\end{center}

\subsection{Sizing Rule}

The code uses different sizing conventions pre-flop and post-flop:

\begin{itemize}[leftmargin=*]
  \item \textbf{Pre-flop}: bucket by total raise size in big blinds. This avoids
  classifying a normal 2.25bb open as ``large'' merely because the blind pot is
  only 1.5bb.
  \item \textbf{Post-flop}: bucket by contribution divided by pot before action.
\end{itemize}

The current pre-flop thresholds are:

\[
\begin{cases}
\text{small} & \text{if raise size} \leq 3\text{bb},\\
\text{medium} & \text{if } 3\text{bb} < \text{raise size} \leq 8\text{bb},\\
\text{large} & \text{if raise size} > 8\text{bb}.
\end{cases}
\]

The current post-flop thresholds are:

\[
\text{pot fraction}
=
\frac{\text{new chips contributed}}{\text{pot before action}},
\]

\[
\begin{cases}
\text{small} & \text{if fraction} < 0.5,\\
\text{medium} & \text{if } 0.5 \leq \text{fraction} < 1.0,\\
\text{large} & \text{if fraction} \geq 1.0.
\end{cases}
\]

\subsection{Key Data Class}

\texttt{ActionEncoding} stores:

\begin{itemize}[leftmargin=*]
  \item \texttt{bucket}: model action bucket id;
  \item \texttt{label}: readable action label;
  \item \texttt{amount\_to}: total amount committed to on the street;
  \item \texttt{contribution}: new chips added by the actor;
  \item \texttt{pot\_fraction}: contribution divided by pot before action;
  \item \texttt{legacy\_bucket}: original parser bucket.
\end{itemize}

\subsection{Important Functions}

\begin{description}
  \item[\texttt{committed\_before\_action(state, hand, street, player)}]
  Reconstructs how much the player has already committed this street.
  It handles pre-flop blinds explicitly because blinds are not stored in the
  betting history.
  \item[\texttt{encode\_action(state, hand, street, legacy\_action)}]
  Produces an \texttt{ActionEncoding} from a pre-action state and legacy action.
\end{description}

\section{\texttt{state\_encoder.py}}

\subsection{Role}

\texttt{state\_encoder.py} compresses the rich parser state into a small
discrete \texttt{StateKey}. This is essential. EM cannot work if every exact
betting history is treated as a different state. The state encoder deliberately
uses coarse buckets.

\subsection{\texttt{StateKey}}

\texttt{StateKey} contains seven fields:

\begin{center}
\begin{tabular}{ll}
\toprule
\textbf{Field} & \textbf{Meaning} \\
\midrule
\texttt{street} & pre-flop, flop, turn, or river \\
\texttt{position} & coarse position bucket \\
\texttt{active\_players} & heads-up, three-way, or multiway \\
\texttt{facing\_bet} & whether the player must call a bet \\
\texttt{raise\_count} & 0, 1, or 2+ raises this street \\
\texttt{spr} & stack-to-pot ratio bucket \\
\texttt{board\_texture} & no board, dry, paired, monotone, connected, etc. \\
\bottomrule
\end{tabular}
\end{center}

The string form is:

\begin{lstlisting}
street|position|active_players|facing_bet|raise_count|spr|board_texture
\end{lstlisting}

\subsection{Important Functions}

\begin{description}
  \item[\texttt{position\_bucket(player, num\_players)}] Uses the parser's seat
  assumptions: \texttt{p1} is small blind and \texttt{p2} is big blind.
  \item[\texttt{active\_count\_bucket(count)}] Buckets active player count into
  \texttt{heads\_up}, \texttt{three\_way}, or \texttt{multiway}.
  \item[\texttt{raise\_count\_bucket(state)}] Reads the parser's street action
  level and buckets into \texttt{raises\_0}, \texttt{raises\_1}, or
  \texttt{raises\_2plus}.
  \item[\texttt{stack\_to\_pot\_bucket(state, player)}] Buckets SPR into
  \texttt{spr\_low}, \texttt{spr\_mid}, or \texttt{spr\_high}.
  \item[\texttt{board\_texture\_bucket(board)}] Produces a coarse board texture
  label.
  \item[\texttt{current\_bet\_to\_call(state, hand, street)}] Reconstructs the
  current amount to call before the action.
  \item[\texttt{facing\_bet\_bucket(state, hand, street, player)}] Compares the
  current bet to the player's committed amount.
  \item[\texttt{encode\_state(state, hand, street, player, action\_encoding)}]
  Returns the final \texttt{StateKey}.
\end{description}

\subsection{Design Tradeoff}

The current state key is intentionally simple. It should be good enough to begin
filtering and EM. If the state space is too coarse, the model will miss useful
strategic context. If it is too detailed, likelihood tables will be sparse and
EM will become unstable.

\section{\texttt{records.py}}

\subsection{Role}

\texttt{records.py} is the bridge between parsing and modeling. It converts
nested legacy parser objects into a flat list of \texttt{DecisionRecord} rows.
Every row corresponds to one player decision.

\subsection{\texttt{DecisionRecord}}

The \texttt{DecisionRecord} data class stores:

\begin{longtable}{lp{10cm}}
\toprule
\textbf{Field} & \textbf{Meaning} \\
\midrule
\texttt{session\_id} & Source session directory. \\
\texttt{hand\_id} & Hand id from PHH metadata when available. \\
\texttt{hand\_index} & Index within parsed session. \\
\texttt{player\_id} & Acting player, e.g. \texttt{p4}. \\
\texttt{street} & Street of the decision. \\
\texttt{decision\_index} & Index within that street. \\
\texttt{board} & Public board before the action. \\
\texttt{pot\_before} & Pot before the action. \\
\texttt{stack\_before} & Acting player's stack before the action. \\
\texttt{active\_player\_count} & Number of players still active. \\
\texttt{state\_key} & Structured \texttt{StateKey}. \\
\texttt{state\_key\_str} & String form of \texttt{StateKey}. \\
\texttt{action\_bucket} & Model action bucket. \\
\texttt{action\_label} & Readable action label. \\
\texttt{amount\_to} & Total amount committed to on current street. \\
\texttt{contribution} & New chips added by the action. \\
\texttt{pot\_fraction} & Contribution divided by pot before action. \\
\texttt{legacy\_action\_bucket} & Original parser action bucket. \\
\texttt{hole\_cards} & True hole cards when known. \\
\texttt{hand\_class} & True 169 class when known. \\
\texttt{strength\_bucket} & True post-flop strength bucket when board exists. \\
\bottomrule
\end{longtable}

\subsection{Important Functions}

\begin{description}
  \item[\texttt{iter\_hand\_decision\_records(hand, session\_id, hand\_index)}]
  Yields records from one parsed hand.
  \item[\texttt{extract\_session\_records(session\_path)}] Parses one session
  directory and returns all decision records.
  \item[\texttt{extract\_many\_session\_records(root\_path)}] Extracts records
  from every session under a root directory.
  \item[\texttt{group\_records\_by\_player\_hand(records)}] Groups rows by
  \((\text{session}, \text{hand}, \text{player})\).
  \item[\texttt{preflop\_only(records)}] Filters records to pre-flop actions.
\end{description}

\subsection{Why This Is Necessary}

The parser stores actions and states separately:

\begin{lstlisting}[language=Python]
hand.states[street][t]
hand.actions[street][t]
\end{lstlisting}

The model needs them joined. \texttt{records.py} performs that join and adds the
model encodings. This is the most important infrastructure layer for starting
EM.

\section{\texttt{gto\_prior.py}}

\subsection{Role}

\texttt{gto\_prior.py} implements a simple, transparent approximation to
\[
p_{\mathrm{GTO}}(a \mid h,s).
\]

It is not a real solver output. Its purpose is to provide a reasonable
pre-flop action likelihood interface so filtering can start immediately. Later,
the same interface can be backed by actual solver frequencies or a better
pre-flop chart.

\subsection{Action Probability Interface}

The main class is \texttt{PreflopGTOPrior}. Its most important method is:

\begin{lstlisting}[language=Python]
action_probs(hand_class, state_key) -> dict[int, float]
\end{lstlisting}

It returns probabilities over all five action buckets.

\subsection{Hand Features}

\texttt{HandClassFeatures} stores:

\begin{itemize}[leftmargin=*]
  \item high rank;
  \item low rank;
  \item pair indicator;
  \item suited indicator;
  \item rank gap;
  \item number of broadway cards;
  \item ace indicator;
  \item scalar heuristic strength.
\end{itemize}

The helper \texttt{hand\_class\_features(hand\_class)} computes these features.

\subsection{Heuristic Logic}

The prior uses:

\begin{itemize}[leftmargin=*]
  \item intrinsic hand strength;
  \item position adjustment;
  \item multiway penalty;
  \item prior raise-count penalty;
  \item premium-hand boost;
  \item speculative suited/connected boost.
\end{itemize}

It then builds action scores and applies a softmax. A small probability floor is
mixed into every action. That floor is important: if an action has probability
zero, Bayesian filtering would permanently eliminate some hands.

\subsection{Theta Connection}

This module does not yet implement a learned \(\theta\). However, it is the
natural place where the current model's GTO perturbation will enter. The future
version can use:

\[
\phi_{h,s,\theta}[a]
\propto
p_{\mathrm{GTO}}(a\mid h,s)
\exp\{u_\theta(h,s,a)\}.
\]

Currently, \texttt{aggression} and \texttt{tightness} are simple manual knobs,
not learned latent variables.

\subsection{EM Prior Counts}

\texttt{dirichlet\_alpha(hand\_class, state\_key, kappa)} returns:

\[
\alpha_{h,s,a}
=
\kappa p_{\mathrm{GTO}}(a\mid h,s).
\]

These pseudo-counts are used by \texttt{em\_preflop.py} in the M-step helper.

\section{\texttt{preflop\_filter.py}}

\subsection{Role}

\texttt{preflop\_filter.py} implements Bayesian range filtering over 169
pre-flop classes.

\subsection{Classes}

\begin{description}
  \item[\texttt{FilterStep}] Debug record for one update. It stores the state
  key, observed action, evidence, ESS, top class, and top probability.
  \item[\texttt{PreflopRangeFilter}] Maintains and updates the current range
  \(R_t(h)\).
\end{description}

\subsection{Update}

\texttt{PreflopRangeFilter.update(state\_key, action\_bucket)} performs:

\[
\widetilde{R}_t(h)
=
R_{t-1}(h)
\phi_{h,s_t}[a_t],
\]

\[
R_t(h)
=
\frac{\widetilde{R}_t(h)}
{\sum_{h'} \widetilde{R}_t(h')}.
\]

The denominator is stored as \texttt{evidence}; it is the predictive
probability of the observed action under the current range and prior model.

\subsection{Effective Sample Size}

\texttt{effective\_sample\_size(distribution)} computes:

\[
\mathrm{ESS}(R)
=
\frac{1}{\sum_h R(h)^2}.
\]

ESS is a concentration diagnostic. A uniform range over 169 classes has ESS
169. A highly concentrated range has lower ESS.

\subsection{Convenience Functions}

\begin{description}
  \item[\texttt{filter\_preflop\_records(records)}] Runs the filter over one
  grouped player-hand record sequence.
  \item[\texttt{uniform\_class\_distribution()}] Returns a uniform distribution
  over 169 classes for sensitivity checks.
\end{description}

\section{\texttt{em\_preflop.py}}

\subsection{Role}

\texttt{em\_preflop.py} provides the starter pieces needed for pre-flop EM. It
does not run a full multi-iteration training loop yet. Instead, it exposes the
core mathematical operations separately so they can be inspected and tested.

\subsection{E-Step}

\texttt{e\_step\_preflop(records, prior\_model)} groups records by
\((\text{session}, \text{hand}, \text{player})\), runs the pre-flop filter on
each group, and returns:

\[
q_m(h)
=
\Prb(V_m=h \mid a^m_{1:T_m},s^m_{1:T_m}).
\]

In code, the result is:

\begin{lstlisting}[language=Python]
dict[(session_id, hand_id, player_id), dict[hand_class, probability]]
\end{lstlisting}

\subsection{Soft Counts}

\texttt{expected\_action\_counts(records, responsibilities)} computes:

\[
N_{h,s,a}
=
\sum_m
\sum_t
q_m(h)\mathbf{1}\{s_t=s,\ a_t=a\}.
\]

This is the central EM bridge: true hand labels are replaced by posterior
responsibilities.

\subsection{M-Step Helper}

\texttt{m\_step\_likelihoods(counts, prior\_model, kappa, smoothing)} computes:

\[
\phi_{h,s,a}
=
\frac{
N_{h,s,a}
+ \kappa p_{\mathrm{GTO}}(a\mid h,s)
+ \epsilon
}{
\sum_{a'}
\left[
N_{h,s,a'}
+ \kappa p_{\mathrm{GTO}}(a'\mid h,s)
+ \epsilon
\right]
}.
\]

\texttt{kappa} controls how strongly the GTO-inspired prior affects the learned
likelihood. \texttt{smoothing} is additional nonzero mass for numerical
stability.

\subsection{Evaluation Helper}

\texttt{true\_class\_log\_loss(records, responsibilities)} computes average:

\[
-\log q_m(h_m^\star),
\]

where \(h_m^\star\) is the true hand class from Pluribus. This metric is only
available because Pluribus reveals hole cards.

\subsection{What Is Still Missing for Full EM}

To run full EM, we still need:

\begin{itemize}[leftmargin=*]
  \item an iterative training loop;
  \item train/validation splits;
  \item model serialization;
  \item a learned likelihood class that can replace \texttt{PreflopGTOPrior}
  inside the filter after each M-step;
  \item convergence monitoring based on held-out action log-loss or true-class
  log-loss.
\end{itemize}

\section{\texttt{README.md}}

\subsection{Role}

\texttt{README.md} is a short usage guide. It lists the package modules and
shows two examples:

\begin{enumerate}[leftmargin=*]
  \item extracting records and filtering one player-hand;
  \item running the EM starter functions.
\end{enumerate}

\subsection{Minimal Filtering Example}

\begin{lstlisting}[language=Python]
from utils_jin.records import extract_session_records, group_records_by_player_hand
from utils_jin.preflop_filter import filter_preflop_records

records = extract_session_records("pluribus/78", include_postflop=False)
groups = group_records_by_player_hand(records)

key, one_player_hand = next(iter(groups.items()))
filt = filter_preflop_records(one_player_hand)

print(key)
print(filt.top_k(10))
print(filt.steps[-1])
\end{lstlisting}

\subsection{Minimal EM Starter Example}

\begin{lstlisting}[language=Python]
from utils_jin.em_preflop import (
    e_step_preflop,
    expected_action_counts,
    m_step_likelihoods,
)

q = e_step_preflop(records)
counts = expected_action_counts(records, q)
phi = m_step_likelihoods(counts)
\end{lstlisting}

\section{Current Smoke-Test Results}

The implementation was smoke-tested on the Pluribus data. Full pre-flop
extraction produced:

\begin{center}
\begin{tabular}{lr}
\toprule
\textbf{Quantity} & \textbf{Value} \\
\midrule
Pre-flop decision records & 61{,}700 \\
Distinct compact pre-flop states & 76 \\
Known hand-class labels & 61{,}700 \\
Exact combos generated & 1{,}326 \\
Pre-flop classes generated & 169 \\
\bottomrule
\end{tabular}
\end{center}

The action bucket distribution after re-bucketing was:

\begin{center}
\begin{tabular}{lr}
\toprule
\textbf{Bucket} & \textbf{Count} \\
\midrule
0: fold & 44{,}077 \\
1: check/call & 6{,}777 \\
2: small bet/raise & 8{,}293 \\
3: medium bet/raise & 966 \\
4: large bet/raise & 1{,}587 \\
\bottomrule
\end{tabular}
\end{center}

This confirms that all five model action buckets are currently represented.

\section{How the Scripts Fit Together}

\subsection{Filtering Workflow}

\begin{lstlisting}[language=Python]
from utils_jin.records import extract_session_records, group_records_by_player_hand
from utils_jin.preflop_filter import filter_preflop_records

records = extract_session_records("pluribus/78", include_postflop=False)
groups = group_records_by_player_hand(records)

for key, group in groups.items():
    filt = filter_preflop_records(group)
    posterior = filt.range
\end{lstlisting}

This workflow uses:

\begin{itemize}[leftmargin=*]
  \item \texttt{compat.py} indirectly through \texttt{records.py};
  \item \texttt{records.py} for extraction;
  \item \texttt{action\_encoder.py} and \texttt{state\_encoder.py} inside
  extraction;
  \item \texttt{cards.py} for true class labels and initial prior;
  \item \texttt{gto\_prior.py} for action likelihoods;
  \item \texttt{preflop\_filter.py} for posterior updates.
\end{itemize}

\subsection{EM Starter Workflow}

\begin{lstlisting}[language=Python]
from utils_jin.em_preflop import (
    e_step_preflop,
    expected_action_counts,
    m_step_likelihoods,
    true_class_log_loss,
)

q = e_step_preflop(records)
counts = expected_action_counts(records, q)
phi = m_step_likelihoods(counts)
loss = true_class_log_loss(records, q)
\end{lstlisting}

This workflow computes one EM-style E-step and one M-step table. It does not yet
loop until convergence.

\section{Recommended Next Engineering Steps}

\begin{enumerate}[leftmargin=*]
  \item Add a train/test split utility for Pluribus sessions.
  \item Add a likelihood model class that stores learned \(\phi_{h,s,a}\)
  tables and implements the same action-probability interface as
  \texttt{PreflopGTOPrior}.
  \item Implement an actual EM loop:
  \[
  \text{E-step} \rightarrow \text{soft counts} \rightarrow \text{M-step}
  \rightarrow \text{evaluate} \rightarrow \text{repeat}.
  \]
  \item Add held-out action log-loss and true-class log-loss metrics.
  \item Implement the post-flop range-to-strength mapper:
  \[
  R(h) + \text{board} \rightarrow Q(w).
  \]
  \item Add a post-flop strength filter using \(\psi_{w,s,a}\).
  \item Add learned \(\theta\) only after the baseline EM loop works.
\end{enumerate}

\section{Summary}

\texttt{utils\_jin/} now contains the complete utility layer needed to begin
pre-flop filtering and to start implementing EM. The most important design
achievement is separation of concerns:

\begin{itemize}[leftmargin=*]
  \item parser compatibility is isolated in \texttt{compat.py};
  \item raw poker data is flattened in \texttt{records.py};
  \item model state/action spaces are controlled by encoders;
  \item GTO-inspired priors are isolated behind a replaceable interface;
  \item filtering is independent of extraction;
  \item EM utilities are built on top of filtering outputs.
\end{itemize}

This structure lets the project improve one layer at a time without rewriting
the entire pipeline.

\end{document}
